% Proposed insertion after chapters/07-mixed-dougall-identities.tex.
% Requires amsmath/amssymb/amsthm and the existing theorem environments.
% This file does not modify the manuscript. Full proofs are in article.tex.
\section{Cubic invariant descent for odd harmonic tails}
\label{cti:section}
Put $T_s(n)=\zeta(s)-H_n^{(s)}$ for $s\ge3$ and $n\ge0$.
For $a\ge0$ and $d\ge2$, define
\[
 M_{a,d}=\sum_{n\ge0}(2n+1)
   \sum_{\substack{j_1+\cdots+j_d=a\\j_i\ge0}}
       \prod_{i=1}^dT_{2j_i+3}(n).
\]
The inner tuples are ordered. Define finite polynomials $P_{a,d}$ by
$P_{0,0}=1$, $P_{a,0}=0$ for $a>0$, and
\[
 P_{a,d}=\frac2d\sum_{\substack{1\le b\le d\\b\text{ odd}}}
    \sum_{i=0}^a\binom{i+b-1}{i}T_{2i+3b}P_{a-i,d-b}.
\]
\begin{theorem}[Invariant tail coefficients]\label{cti:coefficients}
For $a\ge0,d\ge2$,
\[
 \sum_{n\ge0}(2n+1)P_{a,d}(T_3(n),T_5(n),\ldots)
 =\frac{2a+3d-3}{a+d-1}\binom{a+d-1}{d-1}\zeta(2a+3d-2).
\]
Every series is ordinarily absolutely convergent. Moreover $M_{a,d}$
has a constructive rational linear expression in convergent multiple zeta
values of weight $2a+3d-2$ and depth at most $\max(1,d-2)$.
\end{theorem}
\begin{proof}
Specialize the first mixed resonance to $a=1$ and put
$\sigma=p^2+pq+q^2$, $\tau=pq(p+q)$.
For the roots $r=(p,q,-p-q)$, its Gamma quotient is
\[
 Q_n=\prod_{j=1}^3\frac{\Gamma(n+1+r_j)}{\Gamma(n+1-r_j)}
     =\prod_{k>n}\frac{k^3-\sigma k+\tau}{k^3-\sigma k-\tau}.
\]
On $|\sigma|+|\tau|<1$,
\[
 \log Q_n=2\sum_{i,b\ge0}\frac{\binom{i+2b}{i}}{2b+1}
   T_{2i+6b+3}(n)\sigma^i\tau^{2b+1}.
\]
The normally convergent resonance is
\[
 \sum_n\left\{(n+\tfrac12)(Q_n-1)-\frac\tau{n+1/2}\right\}
 =\tau\left[\psi(\tfrac12)+\frac12\left\{\psi(2)-\sum_j\psi(1+r_j)\right\}\right].
\]
Newton's identity
\[
 \sum_jr_j^k=\sum_{2i+3b=k}
   \frac{k}{i+b}\binom{i+b}{b}(-1)^b\sigma^i\tau^b
\]
and the digamma Taylor series give the asserted coefficient. The recurrence
is obtained by differentiating the exponential with respect to $\tau$.
Its leading term in the number of tail factors is
\[
 P_{a,d}=\frac{2^d}{d!}[\sigma^a]
       \left(\sum_{j\ge0}T_{2j+3}\sigma^j\right)^d+E_{a,d},
\]
where every monomial of $E_{a,d}$ has at most $d-2$ factors: all nonlinear
terms of $\log Q_n$ have odd $\tau$-degree at least three.
Finally, expanding a product of tails and summing $2n+1$ first replaces
the minimum index by its square. Grouping equal indices into ordered
blocks produces a multiple zeta value with the final exponent reduced by
two. It never increases the number of factors. This proves the stated
depth and weight bounds. Normal convergence of the subtracted resonance
justifies coefficient extraction; all nonlinear coefficient series converge
absolutely by $T_s(n)=O(n^{1-s})$.
\end{proof}

\begin{corollary}[Quadratic, cubic and quartic evaluations]\label{cti:moments}
For $a\ge0$,
\[
 M_{a,2}=\frac{2a+3}{2}\zeta(2a+4),\qquad
 M_{a,3}=\frac{(a+1)(2a+7)}4\zeta(2a+7).
\]
In particular,
\[
 \sum_n(2n+1)T_3(n)^3=\frac74\zeta(7),\qquad
 \sum_n(2n+1)T_3(n)^4=\zeta(10)+\zeta(5)^2+2\zeta(7,3).
\]
\end{corollary}
\begin{proof}
The coefficients of $\tau^2$ and $\tau^3$ are
$2A_n^2$ and $4A_n^3/3+2C_n/3$, where
$A_n=\sum_jT_{2j+3}(n)\sigma^j$ and
$C_n=\sum_j\binom{j+2}{2}T_{2j+9}(n)\sigma^j$.
Use $\sum_n(2n+1)T_s(n)=\zeta(s-2)$.
The coefficient $P_{0,4}=2T_3^4/3+4T_3T_9/3$ gives the fourth moment in
terms of $\zeta(10)+\zeta(3,7)+\zeta(9,1)$.
Euler's formula for $\zeta(9,1)$ and the stuffle relation for
$\zeta(3)\zeta(7)$ give the displayed simplification.
\end{proof}

\begin{theorem}[A symmetric digamma cube]\label{cti:digamma}
Set
\[
 A_n(x)=\frac{2\psi(n+1)-\psi(n+1-x)-\psi(n+1+x)}{2x^2},\qquad
 J(x)=\psi(1+x)+\psi(1-x)-2\psi(1).
\]
For $x\in\mathbb C\setminus(\mathbb Z\setminus\{0\})$,
\[
 \sum_{n\ge0}(2n+1)A_n(x)^3
 =\frac{J(x)}{4x^6}+\frac{J'(x)}{8x^5}
       -\frac{J''(x)}{16x^4}+\frac{3\zeta(3)}{4x^4}.
\]
At zero use the removable value. The powers are algebraic.
\end{theorem}
\begin{proof}
For $|x|<1$, $A_n(x)=\sum_jT_{2j+3}(n)x^{2j}$.
Apply the operator $(D+1)(2D+7)/4$, $D=(x/2)\partial_x$, to
$(-J(x)/(2x^2)-\zeta(3)-x^2\zeta(5))/x^4$.
This is precisely the cubic generating family just proved.
The series is locally normally convergent away from nonzero integers,
so meromorphic continuation proves the full statement.
\end{proof}

\paragraph{Scope.}
This is a representation-depth upper bound, not an arithmetic independence
claim. It does not separately reduce every unsymmetrized product, and it
does not resolve the Gaussian $S_6$ or revised $S_8$ conjectures.
The full report also proves all-index cubic Hurwitz/Stieltjes traces,
Lerch order-derivative identities and anchored primitives; those results
can be integrated independently of this insert.
